\section{Threat Model and System Application}
\label{sec:threat}

\subsection{System Application}

The target system is a \textbf{Raspberry Pi Pico} (RP2040, dual-core ARM
Cortex-M0+ at 125\,MHz) executing AES-128 in ECB mode using the
\texttt{tiny-AES-c} library~\cite{tinyaes}.
The device receives 16-byte plaintexts over USB serial, encrypts them, and
returns the ciphertext.
A dedicated GPIO pin (GPIO\,22) is asserted high immediately before
\texttt{AES\_ECB\_encrypt} and low immediately after, providing a hardware
trigger signal to a bench oscilloscope.

\subsection{Adversary Threat Model}

We assume a \textbf{local, non-invasive adversary} with the following
capabilities and constraints:

\begin{description}
  \item[Physical access.] The adversary has stable, uninterrupted physical
    access to the target device and can position a power or EM probe over the
    device's power supply rail or decoupling capacitors for the duration of the
    experiment.

  \item[Chosen-plaintext oracle.] The adversary can supply arbitrary
    plaintexts to the device and observe the corresponding ciphertexts.
    This models a scenario where the attacker controls the input (e.g., an
    attacker-controlled message encrypted under a device's stored key).

  \item[Algorithm knowledge.] The adversary knows the cryptographic algorithm
    (AES-128 ECB) and its implementation (standard S-Box, key schedule).
    The \emph{secret key} is unknown.

  \item[Measurement capability.] The adversary can record power traces at
    sufficient sample rate and resolution.
    In the physical setup, a digital storage oscilloscope samples at
    $\geq$\,250\,MS/s over the trigger window.
\end{description}

\subsection{What the Adversary Does Not Have}

\begin{itemize}
  \item No access to the firmware binary (grey-box model assumed for analysis
        purposes; the S-Box structure is public knowledge for AES).
  \item No ability to modify the firmware or introduce hardware faults
        (fault injection is out of scope).
  \item No network or software access to the device; all interaction is
        through the physical serial interface.
\end{itemize}

\subsection{Success Criterion}

A successful attack recovers the full 128-bit secret key with Guessing
Entropy $\text{GE} = 1$ for all 16 key bytes; that is, the correct key byte
ranks first among all 256 hypotheses after processing $N$ traces.
